% GENERATED FILE. Edit canonical lessons, then run npm run generate:books.
% canonical-source-hash: fnv1a64:f4237657d450726f
% canonical-lessons: GE-C268-augenblick, GE-C268-zeitpunkt, GE-C268-zeitraum, GE-C268-dauer, GE-C268-frist

\chapter{Der Augenblick}
\label{ch:ge-der-augenblick}
\hlchaptermodality{268}

\begin{chapteropening}
\textbf{By the end of this chapter:} \emph{I can talk about moments and stretches of time: a moment, a point in time, a period, the length and a deadline.}

Complete the last lesson of chapter 268: I can talk about moments and stretches of time: a moment, a point in time, a period, the length and a deadline.
\end{chapteropening}

\section[der Augenblick]{der Augenblick — the moment, an instant}

\label{lesson:GE-C268-augenblick}



\begin{quote}
Before the new one: say the German for to continue, then the German for to interrupt.
\end{quote}

\subsection*{You'll want to know: der Augenblick}
\textbf{der Augenblick} — \textquotedblleft{}the moment, an instant\textquotedblright{}.

\begin{quote}
Im letzten Augenblick sprang er in den Zug.
\end{quote}

\begin{quote}
\emph{At the last moment he jumped onto the train.}
\end{quote}

Plural: \emph{die Augenblicke}. \emph{Im Augenblick} — at the moment, right now.

\begin{cousinweb}
\emph{Augenblick} comes from Middle High German \emph{ougenblick} \textquotedblleft{}a quick glance\textquotedblright{} (EtymDB), from \emph{Auge} \textquotedblleft{}eye\textquotedblright{} and \emph{Blick} \textquotedblleft{}look\textquotedblright{}: the time an eye takes to blink.
\end{cousinweb}

\begin{grammarlens}[title={What you've built}]
One more word for telling, explaining and describing at B1.
\end{grammarlens}

\subsection*{Your turn}
\begin{itemize}
\raggedright
  \item \emph{Say it:} \emph{der Augenblick}
  \item \emph{Say it:} \emph{der Augenblick}, once more
  \item \emph{Say it:} a sentence of your own with \emph{der Augenblick}
  \item \emph{Recall:} say the German for to continue, then the German for to interrupt, then say \emph{der Augenblick} again
\end{itemize}

\begin{tcolorbox}[breakable,colback=teal!4,colframe=teal!35!black,title={Before you move on}]
What does der Augenblick mean? (\textquotedblleft{}The moment, an instant\textquotedblright{}.) Say it once more.
\end{tcolorbox}
\section[der Zeitpunkt]{der Zeitpunkt — the point in time, the moment when}

\label{lesson:GE-C268-zeitpunkt}



\begin{quote}
Before the new one: say the German for to interrupt, then the German for the moment.
\end{quote}

\subsection*{You'll want to know: der Zeitpunkt}
\textbf{der Zeitpunkt} — \textquotedblleft{}the point in time, the moment when\textquotedblright{}.

\begin{quote}
Zu diesem Zeitpunkt war ich noch in der Schule.
\end{quote}

\begin{quote}
\emph{At that point in time I was still at school.}
\end{quote}

Plural: \emph{die Zeitpunkte}. \emph{Zum richtigen Zeitpunkt} — at the right moment.

\begin{cousinweb}
\emph{Zeitpunkt} joins the inherited \emph{Zeit} \textquotedblleft{}time\textquotedblright{} and \emph{Punkt} \textquotedblleft{}point\textquotedblright{}, which came from Latin \emph{punctum} \textquotedblleft{}a prick, a point\textquotedblright{} (EtymDB).
\end{cousinweb}

\begin{grammarlens}[title={What you've built}]
One more word for telling, explaining and describing at B1.
\end{grammarlens}

\subsection*{Your turn}
\begin{itemize}
\raggedright
  \item \emph{Say it:} \emph{der Zeitpunkt}
  \item \emph{Say it:} \emph{der Zeitpunkt}, once more
  \item \emph{Say it:} a sentence of your own with \emph{der Zeitpunkt}
  \item \emph{Recall:} say the German for to interrupt, then the German for the moment, then say \emph{der Zeitpunkt} again
\end{itemize}

\begin{tcolorbox}[breakable,colback=teal!4,colframe=teal!35!black,title={Before you move on}]
What does der Zeitpunkt mean? (\textquotedblleft{}The point in time, the moment when\textquotedblright{}.) Say it once more.
\end{tcolorbox}
\section[der Zeitraum]{der Zeitraum — the period, the span of time}

\label{lesson:GE-C268-zeitraum}



\begin{quote}
Before the new one: say the German for the moment, then the German for the point in time.
\end{quote}

\subsection*{You'll want to know: der Zeitraum}
\textbf{der Zeitraum} — \textquotedblleft{}the period, the span of time\textquotedblright{}.

\begin{quote}
In diesem Zeitraum gab es drei Unfälle.
\end{quote}

\begin{quote}
\emph{In this period there were three accidents.}
\end{quote}

Plural: \emph{die Zeiträume}. Common in reports: \emph{im Zeitraum von Mai bis Juli}.

\begin{cousinweb}
\emph{Zeitraum} joins \emph{Zeit} \textquotedblleft{}time\textquotedblright{} and \emph{Raum} \textquotedblleft{}room, space\textquotedblright{}, from Old High German \emph{rūm} (EtymDB): a stretch of room in time. Both are inherited Germanic words.
\end{cousinweb}

\begin{grammarlens}[title={What you've built}]
One more word for telling, explaining and describing at B1.
\end{grammarlens}

\subsection*{Your turn}
\begin{itemize}
\raggedright
  \item \emph{Say it:} \emph{der Zeitraum}
  \item \emph{Say it:} \emph{der Zeitraum}, once more
  \item \emph{Say it:} a sentence of your own with \emph{der Zeitraum}
  \item \emph{Recall:} say the German for the moment, then the German for the point in time, then say \emph{der Zeitraum} again
\end{itemize}

\begin{tcolorbox}[breakable,colback=teal!4,colframe=teal!35!black,title={Before you move on}]
What does der Zeitraum mean? (\textquotedblleft{}The period, the span of time\textquotedblright{}.) Say it once more.
\end{tcolorbox}
\section[die Dauer]{die Dauer — the length, the duration}

\label{lesson:GE-C268-dauer}



\begin{quote}
Before the new one: say the German for the point in time, then the German for the period.
\end{quote}

\subsection*{You'll want to know: die Dauer}
\textbf{die Dauer} — \textquotedblleft{}the length, the duration\textquotedblright{}.

\begin{quote}
Für die Dauer der Reise blieb das Handy aus.
\end{quote}

\begin{quote}
\emph{For the length of the journey the phone stayed off.}
\end{quote}

No plural. \emph{Auf Dauer} — in the long run.

\begin{cousinweb}
\emph{Dauer} is built on \emph{dauern} \textquotedblleft{}to last\textquotedblright{}, which goes back through Middle High German \emph{dūren} to Latin \emph{dūrāre} \textquotedblleft{}to last\textquotedblright{} (EtymDB).
\end{cousinweb}

\begin{grammarlens}[title={What you've built}]
One more word for telling, explaining and describing at B1.
\end{grammarlens}

\subsection*{Your turn}
\begin{itemize}
\raggedright
  \item \emph{Say it:} \emph{die Dauer}
  \item \emph{Say it:} \emph{die Dauer}, once more
  \item \emph{Say it:} a sentence of your own with \emph{die Dauer}
  \item \emph{Recall:} say the German for the point in time, then the German for the period, then say \emph{die Dauer} again
\end{itemize}

\begin{tcolorbox}[breakable,colback=teal!4,colframe=teal!35!black,title={Before you move on}]
What does die Dauer mean? (\textquotedblleft{}The length, the duration\textquotedblright{}.) Say it once more.
\end{tcolorbox}
\section[die Frist]{die Frist — the deadline, the time limit}

\label{lesson:GE-C268-frist}



\begin{quote}
Before the new one: say the German for the period, then the German for the length.
\end{quote}

\subsection*{You'll want to know: die Frist}
\textbf{die Frist} — \textquotedblleft{}the deadline, the time limit\textquotedblright{}.

\begin{quote}
Die Frist endet am Freitag.
\end{quote}

\begin{quote}
\emph{The deadline ends on Friday.}
\end{quote}

Plural: \emph{die Fristen}. \emph{Innerhalb der Frist} — within the time limit.

\begin{cousinweb}
\emph{Frist} is inherited from Old High German \emph{frist}, from Proto-Germanic \emph{fristiz} (EtymDB), and is a cousin of the old English word \emph{frist} \textquotedblleft{}a period of grace\textquotedblright{}.
\end{cousinweb}

\begin{grammarlens}[title={What you've built}]
One more word for telling, explaining and describing at B1.
\end{grammarlens}

\subsection*{Your turn}
\begin{itemize}
\raggedright
  \item \emph{Say it:} \emph{die Frist}
  \item \emph{Say it:} \emph{die Frist}, once more
  \item \emph{Say it:} a sentence of your own with \emph{die Frist}
  \item \emph{Recall:} say the German for the period, then the German for the length, then say \emph{die Frist} again
\end{itemize}

\begin{tcolorbox}[breakable,colback=teal!4,colframe=teal!35!black,title={Before you move on}]
What does die Frist mean? (\textquotedblleft{}The deadline, the time limit\textquotedblright{}.) Say it once more.
\end{tcolorbox}
